\begin{textsample}{POS Dim 1 – ce – Score 45.00 – ce/ce\_inf\_7.txt}  \label{ex:f1_pos_050}
4.1.1.5.
Aprender a partir das interações e da brincadeira O processo educativo, \textbf{que} visa a possibilitar condições adequadas de educação e cuidados, requer a garantia do respeito ao ritmo próprio de cada criança, além da compreensão de \textbf{que} ela aprende desde o nascimento, nas experiências cotidianas \textbf{vividas}, por meio do corpo, da emoção, da linguagem verbal, das interações e das múltiplas linguagens.
Aprender pode ser entendido como o processo de modificação do modo de agir, sentir e pensar de cada pessoa, e \textbf{que} não pode ser atribuído à \textbf{mera} maturação orgânica, mas à construção de sentidos e significados sobre si mesmo, sobre os outros e sobre o mundo \textbf{que} se dá pela multiplicidade das experiências \textbf{vividas}.
Nessa concepção, as possibilidades de aprendizagem não são resultado de processos espontâneos, mas requerem alguns elementos mediadores, em especial, a colaboração de diferentes parceiros na realização de alguma \textbf{tarefa}.
O ensinar pode ser compreendido como \textbf{um} processo dinâmico, interativo, de \textbf{inúmeras} possibilidades \textbf{que} contribuem, por meio de experiências inaugurais ou ressignificadas, à inserção dos \textbf{sujeitos} ao meio sociocultural.
Esse processo se concretiza em situações diversas e cotidianas.
Por exemplo, quando se \textbf{diz} para \textbf{um} bebê \textbf{que} é hora de mamar e \textbf{que} ele deve estar com fome, ou quando é chamado pelo seu nome, ou ainda quando escuta o nome das pessoas \textbf{que} circulam em seu meio.
Esses significados culturais poderão ser apropriados a partir das experiências desses \textbf{sujeitos}, por isso se ressalta \textbf{que} é na interação com o outro \textbf{que} se dá significado às coisas.
Fica mais fácil compreender essa ideia quando se pensa em situações \textbf{que} são \textbf{vividas} pelas pessoas e quando cada \textbf{uma} relata o fato interpretado, a partir do seu ponto de vista.
Essa perspectiva para a Educação Infantil requer aportes \textbf{teóricos} e práticos, consistentemente fundamentados nas necessidades e características de crescimento e desenvolvimento de bebês, crianças bem pequenas e crianças pequenas.
A presença de professoras, de professores \textbf{que} se portam de forma atenta, sensível, respeitosa, mediadora das valiosas e significativas aprendizagens das crianças, \textbf{que} \textbf{instigam} seus afetos, suas sensações, seus movimentos, possibilitando as interações e as brincadeiras, sua memória, linguagem e identidade, ampliarão as possibilidades de construção de conhecimentos já \textbf{sistematizados} na cultura, e, especialmente, dos sentidos pessoais \textbf{que} decorrem das interpretações infantis.
Mesmo quando \textbf{uma} professora, \textbf{um} professor propõe \textbf{que} as crianças se envolvam em \textbf{uma} atividade em princípio similar, como por exemplo, o desenho ; se houver liberdade de ação e interação pelas crianças na criação desse desenho — o \textbf{que} é \textbf{demandado} quando se lê com atenção as DCNEI e a BNCC — cada criança produzirá graficamente o suporte com \textbf{singularidade}, lançando mão de suas experiências e conhecimentos prévios, interagindo com aquela situação, constituída naquele momento.
Algumas crianças fazem o desenho em silêncio, outras precisam conversar com o colega do \textbf{lado}, outras se levantam da sua cadeira e ficam em pé para experimentar diferentes possibilidades de produção.
Enfim, a proposta de desenhar pode até ser feita coletivamente, mas as experiências das crianças são as mais diversas.
Essa diversidade só enriquece o ambiente de aprendizagem.
Dessa forma, ações de ensino e aprendizagem podem partir de adultos, crianças e situações interativas do cotidiano, mas não se esgotam nesses elementos.
Em outras palavras, não se aprende só com a professora, o professor, mas com diferentes elementos simbólicos \textbf{que} agem como recursos na relação da criança com o mundo. - amplia o olhar para as diferentes fontes de aprendizagem ( adultos, crianças e situações diversificadas de interação ) ; - opõe -se à ideia de ensino como movimento \textbf{que} parte apenas da professora, do professor e \textbf{que} toma a criança como \textbf{mero} receptor de suas mensagens ; - considera \textbf{que} o processo de ensino e aprendizagem \textbf{depende} da interação \textbf{que} se estabelece entre adultos e crianças, e suas decorrências \textbf{dependerão} basicamente da atividade de cada criança, \textbf{que} continuamente atribui sentidos aos significados \textbf{que} lhe são apresentados, sem \textbf{que} esse reconhecimento enfraqueça a importância das ações da professora, do professor.
Com base nessa concepção, propõe -se pensar qual é o papel do professor.
Para tanto, tem -se \textbf{que} afastar a ideia da professora, do professor, tradicionalmente, associado a alguém \textbf{que} é \textbf{um} transmissor de conhecimentos às crianças.
O desafio é buscar \textbf{uma} nova forma de pensar como o professor deve atuar junto às/com as crianças pequenas, bem pequenas e junto aos/com os bebês em articulação com as famílias e os demais profissionais da escola de educação infantil.
No caso da instituição educacional, o professor é \textbf{um} mediador das aprendizagens das crianças.
Ele atua de modo : - indireto, pelo arranjo do ambiente de aprendizagem das crianças, onde outros mediadores estão presentes : os espaços, os objetos, as indumentárias, os livros, os horários, os agrupamentos infantis ; - direto, quando, por exemplo, ao interagir com as crianças, lhes apresenta \textbf{novidades}, dando suporte para elaboração, validação ou negação das hipóteses ; quando as pega no colo numa situação de insegurança ; quando medeia normas de conduta, e outras ações.
Além das interações dos adultos com as crianças, a BNCC ( BRASIL, 2017 ) ressalta também as interações das crianças com outras crianças, \textbf{que} são parceiras na fascinante \textbf{tarefa} de compreender o mundo e a si mesmas.
As crianças, nas interações \textbf{que} estabelecem entre si, aprendem a : - fazer amigos, negociar significados e decisões, resolver conflitos, partilhar sentimentos e combater estereótipos e preconceitos \textbf{que} limitam o desenvolvimento de \textbf{uma} pessoa ; - \textbf{viver} em grupo, a ser sensível ao ponto de vista ou aos sentimentos do outro, a cooperar em diferentes \textbf{tarefas}, a conhecer suas limitações e possibilidades, a aceitar -se e aos companheiros, a controlar seus impulsos e emoções, e a desenvolver variadas formas de comunicação e de expressão de afetos ; - partilhar com outras crianças os conhecimentos e a identidade \textbf{que} o grupo lhes oferece ; - construir sua imagem pessoal no grupo.
Tanto as interações da/do professora/professor com as crianças, das crianças entre si quanto às interações com os espaços e materiais estimulam processos de aprendizagem \textbf{que} fazem avançar o desenvolvimento.
Em \textbf{uma} situação, as explicações dadas pela/pelo professora/professor, ou a narrativa de \textbf{um} caso por outra criança, possibilitam a aprendizagem.
Em outra ocasião, a escuta de \textbf{uma} história, o folheio de \textbf{um} livro, a participação em \textbf{um} faz de conta, a construção de \textbf{um} castelo com sucata, e outras atividades, são \textbf{poderosos} mediadores da conquista pela criança de novas formas de agir, pensar e sentir.
As interações das crianças As interações sociais constituem meios para os processos de aprendizagem e desenvolvimento humano, já \textbf{que}, em seu acontecimento, colocam em jogo produções e ferramentas da cultura.
É por essa importância \textbf{que} elas são consideradas como \textbf{um} dos eixos do trabalho pedagógico dedicado aos bebês e crianças, nos espaços educacionais coletivos.
Na Educação Infantil, sob a \textbf{ótica} das crianças, essas interações podem ocorrer entre : - As crianças e as/os professoras/professores adultas/adultos — essenciais para dar riqueza e complexidade às brincadeiras e demais ações promovidas no cotidiano, seja na instituição de educação infantil e em diferentes contextos de socialização da criança ; - Crianças e crianças — permite o compartilhamento de experiências e informações entre as crianças, propiciando ampliarem ou construírem conhecimentos e se desenvolverem na relação eu — outro.
Portanto, as interações sociais ocorrem nas relações das crianças com o adulto, da criança com a criança, da criança com a cultura, por meio dos processos de mediação simbólica, existentes nos objetos, nas expressões, gestos, na brincadeira, nos ambientes, na palavra, e nas manifestações das múltiplas linguagens da criança ( ARAÚJO, 2017 ).
Especificamente nas instituições de Educação Infantil, é esperado \textbf{que} as interações ocorram em situações concretas e cotidianas, significativas, transformadoras dos espaços em ambientes vivos e dinâmicos, nos quais as crianças se \textbf{revelam} como protagonistas, questionam, dialogam, convivem e participam, construindo sua identidade individual e coletiva.
Tal possibilidade diverge de \textbf{uma} perspectiva na qual a criança só aprende ao escutar atentamente o adulto.
Ao proporcionar diversas experiências, sejam livres ou dirigidas e desafiantes, tais ações contribuem para \textbf{que} as crianças expressem de diversas maneiras seus sentimentos, argumentos, hipóteses, saberes e atuem ativamente na cultura e, com os seus parceiros, produzam as culturas infantis.
Vale ressaltar ainda \textbf{que} existem situações em \textbf{que} “ [... ] as interações professora/professor-criança são mais necessárias ” \textbf{uma} vez \textbf{que} este é profissional de referência para a criança, apoiando as iniciativas infantis, acolhendo os medos, as conquistas, as curiosidades, as alegrias e questionamentos.
A interação qualitativa \textbf{implica} não apenas promover boas experiências para as crianças, mas significa acolhê -las desde \textbf{que} ingressem na instituição, promover o seu bem-estar, considerar o \textbf{que} elas \textbf{dizem}, “ [... ] perceber e respeitar suas características, compreender seus motivos e estabelecer vínculos afetivos [... ] ” ( SEDUC, 2011, p. 38 ). \textbf{Enquanto} em outras situações, as interações das crianças são os recursos mais importantes para as aprendizagens, especialmente durante as brincadeiras, ocasiões em \textbf{que} realizam imitações, opõem -se, disputam objetos e brinquedos, combinam ações, criam enredos e colocam em jogo novas ideias, improvisam falas, entre outras.
Para \textbf{que} as \textbf{partilhas} entre as crianças ocorram de forma democrática, com segurança afetiva, é necessário \textbf{que} o professor planeje as ações, tendo a participação das crianças nas decisões sobre os espaços, sobre os materiais pedagógicos, livros, brinquedos etc., por meio de experiências \textbf{que} possibilitem a escuta e expressão das múltiplas linguagens.
Por meio de práticas interativas e significativas \textbf{que} \textbf{fomentem} a democracia e a participação, garantindo às crianças vivenciar os princípios éticos, políticos e estéticos, a/o professora/professor de Educação Infantil promoverá o respeito aos direitos de cada criança, o \textbf{que} difere “ [... ] das formas de trabalho \textbf{que} buscam disciplinar as crianças por castigos, ameaças e outras formas de violação dos direitos infantis à proteção, ao cuidado e à educação. ” ( Ibidem ).
Na BNCC, as interações das crianças com as outras crianças e com os adultos são eixo fundamental da organização do trabalho pedagógico.
Os Direitos de Aprendizagem e os Campos de Experiências — e em cada Campo, os Objetivos de Aprendizagem e Desenvolvimento — estão sustentados na ideia de \textbf{que} a interação promove novas e diversificadas aprendizagens.
A brincadeira como atividade da criança É por meio da brincadeira \textbf{que} as crianças constituem a si mesmas, dão significado ao mundo ao seu redor, expressam seus desejos, medos, experimentam papéis, interagem com seus pares e compartilham suas culturas.
O brincar é \textbf{inerente} à cultura das crianças e a Educação Infantil possui \textbf{um} importante papel ao garantir esse direito e ampliar esse repertório \textbf{que} elas já possuem, introduzindo -as nos mais variados tipos de brincadeiras.
Parecer da CNE/CP nº 15/2017, \textbf{que} cita a Resolução ( CNE/CEB nº 5/2009 ), assegura o brincar como \textbf{um} dos direitos de aprendizagem e desenvolvimento da criança, no âmbito da Educação Infantil, ressaltando \textbf{que} é direito da criança : brincar cotidianamente de diversas formas, em diferentes espaços e tempos, com diferentes parceiros ( crianças e adultos ), ampliando e diversificando seu acesso a produções culturais, seus conhecimentos, sua imaginação, sua criatividade, suas experiências emocionais, corporais, sensoriais, \textbf{expressivas}, cognitivas, sociais e relacionais ( INCISO II, p. 53 ).
É importante destacar \textbf{que}, ao garantir o BRINCAR em sua perspectiva mais ampla e diversificada, o projeto pedagógico também caminha para a promoção dos demais Direitos de Aprendizagem e Desenvolvimento.
Afinal, ao BRINCAR as crianças aprendem a CONVIVER ; a PARTICIPAR ativamente em seus grupos de parceiros ; a EXPLORAR linguagens, materiais, espaços, elementos da cultura ; a EXPRESSAR emoções, sentimentos, criações, ideias, hipóteses etc. e, por fim, a CONHECER -SE, já \textbf{que} na interação com o outro, na elaboração das brincadeiras e realização dos jogos, compreende as suas necessidades, ideias, afetos, ao compartilhar afetos e ideias diversas \textbf{que} a colocam em outra perspectiva.
Por essa abrangência, a brincadeira é considerada como \textbf{um} dos eixos do trabalho pedagógico na Educação Infantil.
Promover situações na Educação Infantil por meio de práticas pedagógicas \textbf{que} permitam \textbf{que} as crianças brinquem sozinhas, com seus pares, brincadeiras livres ou dirigidas, é \textbf{uma} forma de garantir sua cidadania.
É importante \textbf{salientar} \textbf{que} a criança não nasce sabendo brincar, assim, essa ação é aprendida por meio das interações crianças/adultos, criança/criança, ou no contato com diferentes materiais e artefatos ( brinquedos, livros, materiais não estruturados, dentre outros ).
Na observação de outras crianças ou quando a professora propõe brincadeiras compostas por regras, as crianças vão aprendendo a criar e a reproduzir a ação do brincar.
Ao participar de brincadeiras com adultos e com outras crianças mais experientes, ela vai aprendendo e se apropriando das regras, \textbf{enquanto} brinca.
Basta \textbf{ver} como as crianças brincam no início do ano e o quanto suas brincadeiras vão ganhando novos elementos no decorrer da convivência e das diversas experiências \textbf{que} \textbf{vive} na Educação Infantil.
Para isso, o trabalho pedagógico deve ser rico em linguagens e manifestações artístico-culturais e em situações de observação do mundo e das relações sociais.
No jogo simbólico da brincadeira, as crianças apropriam -se de informações da cultura do mundo adulto, reportando -se a situações vivenciadas por elas em seus contextos sociais, no entanto não se limitam a interiorizá -la, mas, para além disso, interpretam -na criativamente num processo denominado pelo sociólogo William Corsaro ( 2011 ) de “ reprodução interpretativa ”, \textbf{uma} vez \textbf{que} ao reproduzir \textbf{uma} ação vivenciada, a criança não faz \textbf{uma} cópia daquilo \textbf{que} vivenciou, mas a interpreta a partir da sua leitura do \textbf{vivido} na busca por lidar com inquietações próprias e exclusivas, experimentadas nos modos de ser e de se relacionar com as outras crianças.
Nesse sentido, é \textbf{interessante} enxergar em \textbf{uma} brincadeira de papéis registrada em \textbf{uma} escola de educação infantil, localizada em \textbf{uma} comunidade \textbf{ribeirinha} da Amazônia, os elementos da cultura, do modo de \textbf{viver} da comunidade, compondo a brincadeira das crianças.
Episódio : TOMANDO BANHO NA MARÉ⁶ Participantes : Clara ( 5,013 ), Lucas ( 5 ; 2 ), Marcos ( 4 ; 1 ) e Paulo ( 5,0 ) Lucas, Clara e Marcos brincam em \textbf{um} canto da sala.
Clara está com \textbf{uma} boneca na mão e os meninos com alguns brinquedos, como barcos e peças de \textbf{um} jogo de montar.
Clara inicia \textbf{um} diálogo com os meninos.
Clara — Meu filho já sabe subir na árvore.
Ele sobe e não cai.
Lucas — Não cai?
Então ele já sabe.
Mas ele ainda não sabe tomar banho na maré.
Clara — Sabe.
Lucas — Então, vamos lá no igarapé \textbf{ver} se ele sabe.
Clara — Vamos.
Clara — ( Levanta -se com a boneca na mão e acompanha Lucas, \textbf{que} se dirige a \textbf{um} outro canto da sala.
Marcos e Paulo também o acompanham ).
Marcos — Ele vai é morrer afogado.
Clara — Não vai não.
Ele já sabe nadar.
Marcos — Ela \textbf{diz} \textbf{que} o filho dela sabe tudo. ( As crianças chegam no local proposto por Lucas e ele fala ).
Lucas — Isso \textbf{que} eu quero \textbf{ver}.
Pronto, põe ele na água.
Vamos \textbf{ver} se ele sabe tomar banho na maré. ( Clara pega a boneca e vai colocando bem devagar no chão, como se a estivesse colocando na água ).
Lucas — Agora vem a água da maré.
Olha a água chegando.
Está ficando \textbf{forte}, \textbf{forte}. ( Todos olham para a boneca imóvel no chão ).
Lucas — É.
O filho dela parece \textbf{que} não tem medo.
Acho \textbf{que} ele já sabe tomar banho na maré.
Em \textbf{um} breve episódio de brincadeira simbólica, é possível \textbf{ver} todo o empreendimento das crianças na retomada de \textbf{uma} situação do cotidiano e recriação no espaço da escola.
Pensamento, fala, memória são movimentados pelas crianças \textbf{que} precisam reconstruir \textbf{uma} \textbf{cena} de beira de rio e maré, \textbf{que} está presente em imagem mental \textbf{que} é colocada em ação.
Gestos entram em \textbf{cena} com toda a expressividade das crianças, \textbf{que} assumem papéis e precisam ser convincentes \textbf{umas} com as outras e \textbf{consigo} mesmas.
Atividades da criança \textbf{que}, na brincadeira, a impulsionam a ir além do seu desenvolvimento, do seu cotidiano.
Percebem como as mães da comunidade lidam com seus filhos pequenos e valorizam determinados aprendizados.
Os outros colegas a desafiam como “ bons vizinhos ” \textbf{que} são!
E assim, a brincadeira segue reproduzindo significados da cultura e criando novos sentidos para a brincadeira das crianças.
Situações como essa acontecem cotidianamente na Educação Infantil.
Basta parar e registrar momentos de brincadeira das crianças nos diferentes tempos de sua jornada.
Será fácil perceber o \textbf{que} mais chama a atenção delas nas práticas da cultura \textbf{que} constituem a amplitude do estado do Ceará e a diversidade dos municípios cearenses.
A liberdade de escolha e o protagonismo da criança são essenciais na brincadeira.
O caráter de liberdade se expressa pelo fato de a criança poder escolher, mediante a sua \textbf{vontade}, o tema da brincadeira com a qual quer se envolver e o papel \textbf{que} \textbf{desempenhará} dentro dela, entregando -se com toda a sua emoção e subjetividade na atividade.
Durante a brincadeira a criança construirá o seu universo particular e terá a possibilidade de criar \textbf{uma} situação imaginária, na qual poderá assumir diferentes papéis e criar roteiros : como quando brincam de andar de ônibus e percebemos a \textbf{divisão} de papéis distribuídos e acordados pelas próprias crianças : o motorista, o trocador e os passageiros, cada qual com sua função ; e confrontar -se com dificuldades e angústias de forma criativa ; buscando soluções para resolvê -las ou mesmo liquidá -las em \textbf{um} espaço-tempo, onde podem experimentar o mundo e construir \textbf{uma} compreensão própria sobre as pessoas, os sentimentos e os diversos conhecimentos.
A situação imaginária é \textbf{uma} característica presente na brincadeira das crianças e vai sendo apreendida através da interação e \textbf{estímulos} ambientais.
Por exemplo, antes de aprender a brincar, \textbf{uma} espiga de milho é para a criança apenas \textbf{um} alimento e como tal é apreendido por ela, mas quando a imaginação passa a ocupar \textbf{um} papel importante na conduta infantil, \textbf{uma} espiga de milho pode vir a ser \textbf{uma} boneca e a criança se comporta em relação a ela de acordo com o novo sentido \textbf{que} ela lhe dá. ( VASCONCELOS, 2003, p. 10 ) O brincar também pode envolver a reiteração e a criança demonstra isso repetindo a situação até \textbf{que} se satisfaça, ou se sinta confiante.
Isso pode ocorrer também durante outras atividades prazerosas, como na leitura de \textbf{uma} história, por exemplo, quando solicitam \textbf{que} esta seja contada e recontada várias vezes.
Nesse sentido, Moyles ( 2002 ) afirma \textbf{que} o ato de brincar : > [... ] oferece situações em \textbf{que} as habilidades podem ser praticadas, tanto as físicas quanto as mentais, e repetidas tantas vezes quanto for necessário para a confiança e o domínio.
Além, disso, ele permite a oportunidade de explorar os próprios potenciais e limitações. ( p. 22 ) Repetir brincadeiras e jogos nem sempre estará ligado ao prazer, podendo estar também relacionada à necessidade \textbf{que} a criança tem em compreender algo novo ou angustiante \textbf{que} será repetido até \textbf{que} ela compreenda e \textbf{assimile} esta situação.
Ainda em relação ao aspecto da satisfação, a criança muitas vezes abre mão de seus desejos mais \textbf{imediatos} para continuar participando da brincadeira.
Basta observar no processo de negociação \textbf{que} acontece frequentemente entre as crianças, quando \textbf{uma} delas sugere algo \textbf{que} o grupo recusa e a criança tranquilamente \textbf{diz} : “ Tá bom! ” e continua a brincar.
O prazer, então, é alcançado na manutenção da atividade de brincar e não na realização \textbf{imediata} de desejos.
O jogo deve estar presente no conjunto das situações de aprendizagem e sempre \textbf{que} as crianças mostrem interesse em brincar com seus pares, o professor tem \textbf{uma} excelente oportunidade para observar e registrar como elas se organizam no grupo, como brincam, ou mesmo para observar \textbf{uma} criança \textbf{que} esteja lhe chamando a atenção.
São muitas as oportunidades \textbf{que} as crianças podem ter para brincar e aprender.
Algumas possibilidades de aprendizagem serão descritas a seguir : Os bebês menores de 1 ano podem aprender a : - brincar com os professores de cobrir e descobrir o rosto, de procurar e achar objetos escondidos, de esconder -se em algum canto da sala e ser encontrado ; - encaixar peças de madeira, empilhar cubos, puxar objetos etc. ; - envolver -se em troca de objetos com outros bebês, em ações ritmadas como bater as palmas das mãos sobre \textbf{uma} superfície, em entrar e sair de túneis etc. ; - ouvir e dançar, dentro de suas possibilidades, músicas e ritmos diversificados.
Com as crianças de 1 e 2 anos, novas aprendizagens podem ser estimuladas, tais como : - brincar de roda imitando gestos e cantos do professor e dos colegas ; - brincar de esconde-esconde, de jogar bola e de correr, com a supervisão do professor ; - imitar gestos e vocalizações de adultos, crianças ou animais, dentre outras, ou a imitar objetos ( o som do motor de \textbf{um} carro, por exemplo ) ; - brincar com diferentes ritmos e músicas variadas, vivenciando diversas possibilidades de movimentação corporal ( diferentes posições, velocidades, sensações ).
As crianças de 3 anos podem aprender a : - participar de brincadeiras de roda, sem precisar ter o professor como modelo ; - brincar de esconde-esconde e pega-pega, ou jogar bola, com supervisão do professor ; - imitar gestos, posturas e vocalizações de modelos ( adultos, crianças, animais ou personagens de histórias ) na ausência deles ; - utilizar objetos e vestimentas para assumir papéis no faz de conta onde reproduz situações cotidianas ; - reproduzir as ações de \textbf{um} personagem de \textbf{uma} história lida ( imitar o lobo da história, caminhar como os sete anões cantando na floresta ) ; - construir, ajudadas pelo professor, brinquedos com sucatas a partir de modelos, casas ou castelos com areia, sucata, tocos de madeira e outros materiais.
Crianças de 4 anos podem aprender a : - comunicar -se com os companheiros na brincadeira utilizando sons, musicais ou não, ou diferentes formas de gestos e expressões vocais e corporais ; - brincar com a sonoridade de palavras, criando rimas ; - brincar de esconde-esconde, de jogar bola, de pique, de seguir o mestre, de lenço atrás, de caça ao tesouro etc. ; - montar quebra-cabeça ; - dramatizar \textbf{uma} história usando bonecos ou marionetes como \textbf{atores} ; - construir brinquedos, casas e cidades com diferentes materiais, sem necessariamente usar \textbf{um} modelo.
As crianças de 5 anos podem ser apoiadas a aprender, dentre outras coisas, a : - brincar de pula-sela, amarelinha, corda, pega-pega ; - elaborar coletivamente e dramatizar \textbf{um} enredo usando bonecos como \textbf{atores} ; - escolher vestimenta para compor \textbf{um} personagem para si ou para \textbf{um} colega ; - maquiar -se ou a \textbf{um} colega para \textbf{desempenhar} certo papel ; - discutir as intenções dos personagens de \textbf{um} enredo encenado ; - participar de jogos de tabuleiro como : loto, damas, memória, dominó etc. ; - explicar as regras de \textbf{um} jogo para outra criança ; - fazer brinquedos com sucata sem seguir modelo ; - construir casa/castelos de cartas, de cartolina, de panos e outros materiais ; - fazer dobraduras \textbf{simples}, elaborar máscaras, fazer bonecas de pano, ou de espiga de milho ; - construir e empinar pipas com a ajuda de \textbf{um} adulto. ( SEDUC, 2011, p. 41-42 ).
Essas são algumas das brincadeiras \textbf{que} podem ser organizadas diariamente.
Para tanto, a escolha do local e dos materiais \textbf{que} estarão disponíveis é fundamental.
Com relação à brincadeira, o professor precisa : - oferecer \textbf{um} repertório de cantigas, parlendas, adivinhas etc., possibilitando \textbf{que} as crianças vivenciem brincadeiras dançando, cantando, imitando ; - oportunizar situações em \textbf{que} as crianças possam brincar de faz de conta de diferentes formas : sozinhas, com o grupo, de forma livre e orientada pelo professor ; - estimular situações em \textbf{que} as crianças organizem enredos para as dramatizações, roteiros para a produção de danças e musicais e, ainda, planejem a confecção de brinquedos ; - respeitar o tempo e o ritmo das crianças \textbf{enquanto} brincam ; - mediar os conflitos surgidos nas brincadeiras ; - participar das brincadeiras, se solicitado.
A brincadeira é a atividade principal do dia a dia da criança.
Nas brincadeiras, as crianças podem tomar decisões, expressar sentimentos e valores, conhecer a si, aos outros e o mundo, repetir ações prazerosas, partilhar, expressar sua individualidade e identidade por meio de diferentes linguagens, usar o corpo, os sentidos, os movimentos, solucionar problemas e criar.
É possível fazer \textbf{um} rico exercício com a leitura da BNCC se buscarmos enxergar nos diversos Objetivos de Aprendizagem e Desenvolvimento situações diversas de interações das crianças e brincadeira \textbf{que} mobilizam suas aprendizagens e desenvolvimento.
Os Objetivos de Aprendizagem devem ser lidos sempre no contexto dos Campos de Experiências.
Tanto eles são possibilidades de aprendizagem em \textbf{um} Campo, como eles também \textbf{atravessam} outros Campos.
Por exemplo, se tomarmos o exemplo do jogo simbólico do Tomando Banho na Maré, \textbf{veremos} as crianças interagindo e construindo conhecimentos a partir da constituição de campos como : O eu, o outro e o nós ; Corpo, gestos e movimentos e Escuta, fala, pensamento e imaginação, e, a \textbf{depender} da reflexão, inclusive, outros campos.
Refletir, no contexto das interações reais acontecidas na turma de crianças, como os objetivos foram mobilizados no processo também pode enriquecer o olhar da(o) professora(r) para os processos educativos.
Vale \textbf{aqui} \textbf{um} destaque para o cuidado \textbf{que} devemos ter para não ler os objetivos de aprendizagem e desenvolvimento como objetivos pontuais a atingir ou não atingir.
Se olharmos atentamente para esses organizadores na BNCC, eles nos \textbf{dizem} de processos \textbf{que} vão se desenvolvendo continuadamente e de formas diversificadas.
Por exemplo, na brincadeira, quando as crianças usam gestos e movimentos para \textbf{agregar} novos elementos ao enredo, de certa forma elas criam com o corpo formas diversificadas de expressão, de sentimentos, sensações e emoções, elas mobilizam esses processos e seguem aprendendo sobre elas mesmas, sobre os colegas e suas formas de agir e pensar o mundo, assim como sobre o repertório de sua cultura.
São processos experienciados pelas crianças e pelos adultos \textbf{que} interagem com elas.

% matched lemmas: agregar, aqui, assimilar, ator, atravessar, cena, conseguir, demandar, depender, desempenhar, divisão, dizer, enquanto, estímulo, expressivo, fomentar, forte, imediato, implicar, inerente, instigar, interessante, inúmero, lado, mero, novidade, partilha, poderoso, que, revelar, ribeirinho, salientar, simples, singularidade, sistematizar, sujeito, tarefa, teórico, um, uma, ver, viver, vontade, ótica
\end{textsample}
